\documentclass[spanish,letterpaper,oneside]{article}

\def\documenttitle {La Investigación-Acción para Transformar la Práctica Matemática}
\def\documentsubtitle {Fundamentos para la enseñanza y el aprendizaje I · Cohorte SD26}
\def\documentsubject {Actividad 4 · Sesión 7}
\def\documentauthor {Mtro. Martín Jonathan de la Cruz Muñoz}
\def\coursename {Fundamentos para la enseñanza y el aprendizaje I}
\def\coursecode {MAEM-SD26}

\def\universityname {Instituto de Investigación, Innovación y Estudios de Posgrado para la Educación del Estado de Nuevo León}
\def\universityfaculty {Maestría en Aprendizaje y Enseñanza de las Matemáticas}
\def\universitydepartment {Fundamentos para la enseñanza y el aprendizaje I}
\def\universitydepartmentimage {img/iiiepe/marca-iiiepe.png}
\def\universitydepartmentimagecfg {height=1.75cm}
\def\universitylocation {Monterrey, Nuevo León}
\def\studentfullname {Mtro. Martín Jonathan de la Cruz Muñoz}
\def\studentid {00874}
\def\studentemail {martin.delacruz@campus.iiiepe.edu.mx}
\def\programperiod {2do periodo, septiembre--diciembre de 2026 (SD26)}
\def\programmodality {Modalidad presencial}
\def\programshort {MAEM}

\def\authortable {
  \begin{tabular}{ll}
    \textbf{Alumno:} & \studentfullname \\
    Matrícula: & \studentid \\
    Correo institucional: & \studentemail \\
    Unidad: & \coursename \\
    Programa: & \programshort \\
    Periodo: & \programperiod \\
    & \\
    \multicolumn{2}{l}{Fecha de realización: 19 de noviembre de 2026} \\
    \multicolumn{2}{l}{\universitylocation}
  \end{tabular}
}

\input{template}

\usepackage{helvet}
\renewcommand{\familydefault}{\sfdefault}
\usepackage{array}
\usepackage{longtable}
\usepackage{pdflscape}
\usepackage{tikz}
\usepackage{pgfplots}
\pgfplotsset{compat=1.18}
\usepackage[most]{tcolorbox}
\usetikzlibrary{arrows.meta,positioning,calc,fit,backgrounds,shapes.geometric}

\definecolor{iiiepeNavy}{HTML}{163B5C}
\definecolor{iiiepeBlue}{HTML}{1F6F9F}
\definecolor{iiiepeAqua}{HTML}{20B7A8}
\definecolor{iiiepeOrange}{HTML}{F28C28}
\definecolor{iiiepePaper}{HTML}{F6F8F7}
\definecolor{iiiepeInk}{HTML}{1E2A32}
\definecolor{iiiepePale}{HTML}{EEF5F8}
\definecolor{iiiepeMint}{HTML}{E7F6F4}
\definecolor{iiiepeWarm}{HTML}{FFF3E6}
\def\portraittitlecolor {iiiepeNavy}
\setcitestyle{authoryear,open={(},close={)}}
\makeatletter
\renewcommand{\@biblabel}[1]{}
\makeatother
\setlength{\bibhang}{1.27cm}
\setlength{\bibsep}{6pt}
\AtBeginDocument{\renewcommand{\refname}{Referencias}}

\newcommand{\insertcoverwatermark}{%
  \AddToShipoutPictureBG*{%
    \begin{tikzpicture}[remember picture,overlay]
      \node[opacity=0.04,inner sep=0pt] at (current page.center) {%
        \includegraphics[width=0.65\paperwidth]{img/iiiepe/marca-iiiepe.png}%
      };
    \end{tikzpicture}%
  }%
}

\setlength{\parskip}{5pt}\setlength{\parindent}{0pt}
\begin{document}\insertcoverwatermark\templatePortrait\templatePagecfg\templateFinalcfg\fancyfoot{}\fancyfoot[L]{\scriptsize IIIEPE · MAEM · SD26}\fancyfoot[C]{\scriptsize\thepage}\fancyfoot[R]{\scriptsize Actividad 4 · Investigación-Acción}\setstretch{1.3}\raggedbottom
\section{Introducción y postura central}
La transformación sostenida de la enseñanza matemática exige superar la improvisación anecdótica y las reformas prescritas desde fuera del aula, colocando al profesorado como sujeto reflexivo capaz de indagar rigurosamente sobre su propio quehacer. En torno al papel de la investigación-acción como estrategia para transformar la práctica de enseñanza y a partir de los planteamientos de \citet{restrepo2004investigacion}, sostengo la tesis de que la investigación-acción educativa constituye el dispositivo metodológico más pertinente para producir saber pedagógico situado, porque articula en espiral la crítica de la rutina docente, la planificación fundada de mejoras y la lectura honesta de evidencias recogidas en el propio grupo.

Mi postura es que, lejos de oponer la teoría académica a la práctica cotidiana del taller o del aula universitaria, la investigación-acción convierte los obstáculos reales de comprensión matemática en objetos sistemáticos de estudio. A continuación se responden de forma argumentada las cinco preguntas rectoras establecidas en la actividad, se delimita una problemática concreta extraída de mi enseñanza en bachillerato técnico e ingeniería, se desarrolla su abordaje a lo largo de las cinco fases metodológicas y se fundamenta por qué esta perspectiva sostiene toda innovación educativa genuina.

\section{Fundamentos y fases de la investigación-acción educativa}

\subsection*{1. ¿Qué es la investigación-acción y cuál es su finalidad en el ámbito educativo?}
De acuerdo con \citet{restrepo2004investigacion}, la investigación-acción educativa es una modalidad de indagación autorreflexiva, cíclica y participativa emprendida por los propios actores escolares para examinar críticamente su práctica, comprender las causas que limitan el aprendizaje y construir saber pedagógico transformador. Su finalidad no consiste en acumular descripciones neutrales externas al aula, sino en mejorar la enseñanza y las condiciones reales en que aprende el estudiantado mediante un diálogo continuo entre teoría educativa, contexto escolar y acción reflexiva \citep{medina2009didactica,unesco2015replantear}.

\subsection*{2. ¿Por qué considera al docente como investigador de su propia práctica?}
La investigación-acción reconoce al docente como investigador porque es quien habita cotidianamente las tensiones curriculares, escucha de primera mano las concepciones alternativas de sus alumnos y toma decisiones situadas minuto a minuto \citep{restrepo2004investigacion}. Ningún observador externo puede captar por sí solo el significado pedagógico completo de las interacciones en un grupo; cuando el profesorado documenta sus clases, contrasta sus supuestos implícitos y revisa las producciones estudiantiles con criterio sistemático, transita de aplicador pasivo de programas a constructor crítico de conocimiento didáctico \citep{sawyer2008optimizing}.

\subsection*{3. ¿Cuáles son las fases de la investigación-acción y cómo se relacionan entre sí?}
Retomando la tradición metodológica recuperada por \citet{restrepo2004investigacion} —deconstrucción de la práctica, reconstrucción planificada y validación efectiva— y la secuencia operativa trabajada en la sesión (identificación del problema, planificación, acción, observación y reflexión), las fases no operan como pasos burocráticos lineales sino como una espiral helicoidal abierta (Figura~\ref{fig:act4-espiral}). La identificación crítica alimenta el diseño de acciones; la acción en el aula genera registros durante la observación; y la reflexión sobre dichas evidencias da paso a un nuevo ciclo de ajuste cada vez mejor fundamentado \citep{pimienta2012metodologia}.

\begin{figure}[htbp]
\centering
\begin{tikzpicture}[font=\scriptsize, >=Latex]
  \node[draw=iiiepeNavy,fill=iiiepeNavy!12,very thick,circle,minimum size=2.6cm,text width=2.3cm,align=center,font=\scriptsize\bfseries,text=iiiepeNavy] (center) at (0,0) {SABER\\PEDAGÓGICO\\SITUADO\\{\mdseries(Espiral crítica)}};

  \node[draw=iiiepeNavy,fill=iiiepePale,thick,rounded corners=4pt,text width=4.0cm,align=center,inner sep=5pt] (f1) at (0,3.2) {\textbf{1. Identificación / Deconstrucción}\\Delimitar el obstáculo con evidencias reales de aula.};
  \node[draw=iiiepeBlue,fill=iiiepePale,thick,rounded corners=4pt,text width=4.0cm,align=center,inner sep=5pt] (f2) at (4.7,1.1) {\textbf{2. Planificación / Reconstrucción}\\Diseñar estrategias, andamiajes e instrumentos.};
  \node[draw=iiiepeAqua!80!black,fill=iiiepeMint,thick,rounded corners=4pt,text width=4.0cm,align=center,inner sep=5pt] (f3) at (3.0,-2.65) {\textbf{3. Acción en el aula}\\Implementar la mediación en condiciones reales.};
  \node[draw=iiiepeOrange!90!black,fill=iiiepeWarm,thick,rounded corners=4pt,text width=4.0cm,align=center,inner sep=5pt] (f4) at (-3.0,-2.65) {\textbf{4. Observación sistemática}\\Recoger producciones, diarios y registros de diálogo.};
  \node[draw=iiiepeNavy,fill=iiiepePale,thick,rounded corners=4pt,text width=4.0cm,align=center,inner sep=5pt] (f5) at (-4.7,1.1) {\textbf{5. Reflexión evaluativa}\\Interpretar avances y abrir nuevo ciclo mejorado.};

  \draw[->,very thick,iiiepeNavy] (f1.east) to[bend left=18] (f2.north);
  \draw[->,very thick,iiiepeBlue] (f2.south) to[bend left=18] (f3.east);
  \draw[->,very thick,iiiepeAqua!80!black] (f3.west) to[bend left=12] (f4.east);
  \draw[->,very thick,iiiepeOrange!90!black] (f4.west) to[bend left=18] (f5.south);
  \draw[->,very thick,iiiepeNavy] (f5.north) to[bend left=18] node[above left,font=\scriptsize\bfseries] {Nuevo ciclo} (f1.west);
\end{tikzpicture}
\caption{Ciclo helicoidal de las cinco fases de la investigación-acción educativa y su convergencia en el saber pedagógico.}
\label{fig:act4-espiral}
\end{figure}

\subsection*{4. ¿Qué importancia tienen la observación, la reflexión y la recopilación de evidencias?}
Sin observación intencionada ni acopio de evidencias concretas (cuadernos, matrices de seguimiento, registros de discusión oral, soluciones erróneas comentadas), el análisis propio corre el riesgo de reducirse a opiniones subjetivas o a culpar al desinterés estudiantil \citep{restrepo2004investigacion}. Recopilar evidencias permite descubrir regularidades ocultas: por ejemplo, notar si los fallos ocurren al traducir un enunciado físico a símbolos o al interpretar un valor obtenido, orientando la reflexión hacia decisiones didácticas comprobables \citep{eef2021feedback,bostrom2023formative}.

\subsection*{5. ¿Por qué constituye una alternativa para transformar y mejorar la práctica docente?}
A diferencia de las capacitaciones teóricas desconectadas de los planteles reales, la investigación-acción integra diagnóstico, ensayo pedagógico y evaluación continua en el escenario auténtico donde enseña cada profesor \citep{restrepo2004investigacion}. Este proceso rompe la inercia de repetir durante años las mismas secuencias ineficaces, fortalece el juicio profesional autónomo y consolida comunidades colegiadas comprometidas con la mejora educativa \citep{quiroga2022aprendizaje}.

\section{Aplicación a una problemática delimitada de la práctica matemática}
En mi docencia de matemáticas para bachillerato técnico (mantenimiento industrial y mecatrónica) y primeros semestres de ingeniería, he identificado una problemática recurrente: \textbf{el predominio de la manipulación algebraica mecánica en modelos lineales sin interpretación física de las variables, unidades ni límites de validez}. Este hallazgo se fundamenta en evidencias concretas observadas en tareas de taller, evaluaciones diagnósticas y revisiones de cuadernos: numerosos estudiantes aplican correctamente la fórmula $m = (y_2-y_1)/(x_2-x_1)$ o despejan $x$ en $ax + b = c$, pero al preguntarles qué representa el número obtenido en un circuito o una curva de calibración no distinguen la tasa de cambio de la condición inicial ni detectan valores físicamente imposibles \citep{star2015algebra}. En coherencia con la crítica pedagógica, esta dificultad no puede atribuirse únicamente al alumnado; con frecuencia obedece a secuencias docentes que privilegian la rapidez algorítmica por encima de la modelación contextualizada.

El cuadro siguiente sistematiza el desarrollo operativo por aplicar en un ciclo completo de investigación-acción de cinco fases:

{\small\setstretch{1.15}\renewcommand{\arraystretch}{1.35}
\begin{longtable}{|>{\raggedright\arraybackslash}p{0.18\linewidth}|>{\raggedright\arraybackslash}p{0.41\linewidth}|>{\raggedright\arraybackslash}p{0.35\linewidth}|}\hline
\rowcolor{iiiepePale} \textbf{Fase metodológica} & \textbf{Desarrollo situado en la problemática matemática} & \textbf{Evidencias e indicadores de rigor} \\\hline
\endfirsthead
\hline
\rowcolor{iiiepePale} \textbf{Fase metodológica} & \textbf{Desarrollo situado en la problemática matemática} & \textbf{Evidencias e indicadores de rigor} \\\hline
\endhead
1. Identificación del problema & Reconocer que la enseñanza centrada en despejes abstractos provoca respuestas numéricas correctas pero carentes de significado físico, dimensional y contextual. & Análisis de 20--25 producciones diagnósticas previas, registro de errores frecuentes de unidades y diario de clase inicial. \\\hline
2. Planificación & Rediseñar la secuencia de función lineal mediante tres sesiones basadas en calibración e interpretación de datos simulados o experimentales, incorporando tabla--gráfica--ecuación y preguntas escritas obligatorias de interpretación \citep{sawyer2008optimizing}. & Secuencia didáctica escrita, guía de preguntas generadoras y matriz de evaluación formativa compartida desde el inicio. \\\hline
3. Acción & Implementar trabajo en binas y equipos pequeños; antes de ejecutar cálculos algebraicos, cada equipo anticipa el comportamiento esperado y valida las unidades de la pendiente e intersección; puesta en común dialógica sin anticipar la respuesta docente. & Registro cronológico de implementación, fotografías o copias anonimizadas de hojas de trabajo en equipo. \\\hline
4. Observación & Documentar sistemáticamente cómo argumentan los estudiantes sus modelos, dónde persisten confusiones entre relación proporcional y afín, y qué apoyos desbloquean la explicación conceptual \citep{eef2021feedback}. & Diario de campo docente, lista de cotejo de argumentación escrita y muestras comparadas antes/después de la intervención. \\\hline
5. Reflexión & Analizar si las modificaciones ampliaron la capacidad del grupo para justificar unidades, dominio válido y decisiones con el modelo, identificando ajustes necesarios en los andamiajes para el ciclo posterior \citep{restrepo2004investigacion}. & Matriz comparativa de hallazgos, síntesis de obstáculos persistentes y rediseño documentado de la secuencia. \\\hline
\end{longtable}}

\section{Transferencia profesional y fundamento metodológico de la innovación}
La aplicación de este ciclo en el contexto laboral y formativo de nivel medio superior y universitario transforma la planeación en una hipótesis viva sujeta a verificación ética y empírica. Además, involucra al docente en la revisión continua de sus creencias; sin embargo, no basta con enunciar buenas intenciones sin seguimiento ordenado. Por ello, en lugar de adoptar modas didácticas pasajeras, la investigación-acción provee una ruta disciplinada para validar qué estrategias ayudan efectivamente a jóvenes técnicos e ingenieros a pensar matemáticamente \citep{eef2017mathematics}. En consecuencia, es decir, transformar el aula exige coherencia entre diagnóstico, acción y reflexión evaluativa.

\clearpage
\section{Conclusiones: la investigación-acción como base de la innovación}
En conclusión, considero que la investigación-acción educativa constituye el fundamento metodológico indispensable de toda propuesta seria de innovación educativa. Innovar no equivale a introducir cambios vistosos sin seguimiento, sino a modificar intencionalmente las prácticas de enseñanza a partir de problemas diagnosticados con evidencia real, someter dichas modificaciones a observación honesta en el aula y reconstruir el saber pedagógico mediante reflexión crítica compartida \citep{restrepo2004investigacion}.

Por tanto, al integrar el análisis conceptual de las cinco preguntas con un diseño operativo aplicable a la comprensión de modelos lineales, este análisis consolida el puente metodológico entre la reflexión identitaria sobre la propia práctica docente y el diseño de innovaciones pedagógicas orientadas al aprendizaje profundo.

\clearpage
\begin{thebibliography}{99}
\bibitem[Boström y Palm(2023)]{bostrom2023formative} Boström, E., y Palm, T. (2023). Formative assessment practice in mathematics: Teachers' perceptions of prerequisite conditions. \textit{Frontiers in Education}, 8, Artículo 1101192.
\bibitem[Collin y Quigley(2021)]{eef2021feedback} Collin, J., y Quigley, A. (2021). \textit{Teacher feedback to improve pupil learning}. Education Endowment Foundation.
\bibitem[Henderson et~al.(2017)]{eef2017mathematics} Henderson, P., Hodgen, J., Foster, C., y Kuchemann, D. (2017). \textit{Improving mathematics in Key Stages 2 and 3}. Education Endowment Foundation.
\bibitem[Medina Rivilla y Salvador Mata(2009)]{medina2009didactica} Medina Rivilla, A., y Salvador Mata, F. (2009). \textit{Didáctica general} (2.ª ed.). Pearson Educación.
\bibitem[Pimienta Prieto(2012)]{pimienta2012metodologia} Pimienta Prieto, J. H. (2012). \textit{Metodología constructivista: Guía para la planeación docente} (3.ª ed.). Pearson Educación.
\bibitem[Quiroga Aguilar y Lara Órdenes(2022)]{quiroga2022aprendizaje} Quiroga Aguilar, L., y Lara Órdenes, E. (2022). El aprendizaje profundo como herramienta para cambio en la visión de aprendizaje de una cultura escolar. \textit{Revista Educación las Américas}, 12(1). \url{https://doi.org/10.35811/rea.v12i1.201}
\bibitem[Restrepo Gómez(2004)]{restrepo2004investigacion} Restrepo Gómez, B. (2004). La investigación-acción educativa y la construcción de saber pedagógico. \textit{Educación y Educadores}, (7), 45--55.
\bibitem[Sawyer(2008)]{sawyer2008optimizing} Sawyer, R. K. (2008). Optimizing learning: Implications of learning sciences research. OECD Publishing.
\bibitem[Star et~al.(2015)]{star2015algebra} Star, J. R., Caronongan, P., Foegen, A., Furgeson, J., Keating, B., Larson, M. R., Lyskawa, J., McCallum, W. G., Porath, J., y Zbiek, R. M. (2015). \textit{Teaching strategies for improving algebra knowledge in middle and high school students}. Institute of Education Sciences.
\bibitem[UNESCO(2015)]{unesco2015replantear} UNESCO. (2015). \textit{Replantear la educación: ¿Hacia un bien común mundial?} Ediciones UNESCO.
\end{thebibliography}

\end{document}
